\documentclass[aspectratio=169]{ctexbeamer}

\usepackage{amsmath}
\usepackage{graphicx}
\usepackage{booktabs}
\usepackage{tikz}
\usepackage{listings}
\usepackage{xeCJK}

\usetikzlibrary{positioning, shapes.multipart, arrows.meta, bending, decorations.pathreplacing, trees, automata}
\makeatletter
\def\input@path{{/windows/Users/TanghuiLi/Desktop/latexenv/}}
\makeatother
\usetheme{tongji}
\graphicspath{{~/}}

% Define colors for code syntax highlighting
\definecolor{codegreen}{rgb}{0,0.6,0}
\definecolor{codegray}{rgb}{0.5,0.5,0.5}
\definecolor{codepurple}{rgb}{0.58,0,0.82}
\definecolor{backcolour}{rgb}{0.98,0.98,0.95}

\lstdefinestyle{mystyle}{
    backgroundcolor=\color{backcolour},
    commentstyle=\color{codegreen},
    keywordstyle=\color{magenta},
    numberstyle=\tiny\color{codegray},
    stringstyle=\color{codepurple},
    basicstyle=\footnotesize\ttfamily,
    breakatwhitespace=false,
    breaklines=true,
    captionpos=b,
    keepspaces=true,
    numbers=left,
    numbersep=5pt,
    showspaces=false,
    showstringspaces=false,
    showtabs=false,
    tabsize=2
}
\lstset{style=mystyle}

\lstdefinelanguage{RustLike}{
    morekeywords={fn, let, mut, if, else, while, for, in, loop, break, continue, return, i32},
    sensitive=true,
    morecomment=[l]{//},
    morecomment=[s]{/*}{*/},
    morestring=[b]",
}

%Information to be included in the title page:
\title[Compilation Principles]{%
    编译原理大作业 1：词法和语法分析工具
}
\subtitle{——类Rust语言的词法与语法分析器设计与实现}
\author[李堂辉, 高胜寒, 冉子易]{%
    2351078 李堂辉 \\ 2351837 高胜寒 \\ 2352197 冉子易 \\ \small （按学号排序）
}
\institute[Tongji University]{
    同济大学
}
\date{\today}

\begin{document}

% ===== 标题页 =====
\begin{frame}
    \titlepage
\end{frame}

% ===== 目录 =====
\begin{frame}{目录}
    \tableofcontents
\end{frame}

% ============================================================
\section{项目概述}

\begin{frame}{任务要求}
    \begin{block}{大作业目标}
        设计并实现面向\textbf{类Rust语言}的词法分析和语法分析工具
    \end{block}
    \vfill
    \begin{columns}[T]
        \begin{column}{0.5\textwidth}
            \textbf{基本功能}
            \begin{itemize}
                \item 词法分析：识别Token
                \item 语法分析：构建AST
                \item 输出分析结果
            \end{itemize}
        \end{column}
        \begin{column}{0.5\textwidth}
            \textbf{覆盖规则}
            \begin{itemize}
                \item 必需：0.1--5.1（18条）
                \item 扩展：5.2--9.3（16条）
                \item 合计实现 \textbf{34 条}产生式规则
            \end{itemize}
        \end{column}
    \end{columns}
\end{frame}

\begin{frame}{技术选型}
    \begin{columns}[T]
        \begin{column}{0.5\textwidth}
            \begin{block}{实现语言}
                \textbf{Rust}
                \begin{itemize}
                    \item 类型安全，内存安全
                    \item 枚举类型天然适合Token和AST
                    \item 模式匹配简化分析逻辑
                \end{itemize}
            \end{block}
        \end{column}
        \begin{column}{0.5\textwidth}
            \begin{block}{分析方法}
                \begin{itemize}
                    \item \textbf{词法}：手写Cursor扫描器
                    \item \textbf{语法}：递归下降分析器
                    \item 无外部依赖（纯标准库）
                \end{itemize}
            \end{block}
        \end{column}
    \end{columns}
\end{frame}

% ============================================================
\section{词法分析}

\begin{frame}{词法单元分类}
    \begin{table}
    \centering
    \small
    \begin{tabular}{lp{9cm}}
    \toprule
    \textbf{类别} & \textbf{内容} \\
    \midrule
    关键字 & \texttt{i32, let, if, else, while, return, mut, fn, for, in, loop, break, continue} \\
    标识符 & $(\text{字母}|\_)(\text{字母}|\text{数字}|\_)^*$，不与关键字相同 \\
    数值 & $\text{数字}(\text{数字})^*$ \\
    赋值号 & \texttt{=} \\
    算符 & \texttt{+ - * / == > >= < <= != \&} \\
    界符 & \texttt{( ) \{ \} [ ]} \\
    分隔符 & \texttt{; : ,} \\
    特殊符号 & \texttt{-> . ..} \\
    注释 & \texttt{// ...} 和 \texttt{/* ... */} \\
    结束符 & \texttt{\#} \\
    \bottomrule
    \end{tabular}
    \end{table}
\end{frame}

\begin{frame}[shrink]{词法分析器——状态转换}
    \begin{center}
    \begin{tikzpicture}[
        ->,>=Latex, shorten >=1pt, auto,
        state/.style={circle, draw, minimum size=0.8cm, font=\scriptsize},
        accept/.style={circle, draw, double, minimum size=0.8cm, font=\scriptsize}
    ]
        \node[state] (S) {$S_0$};
        \node[accept] (ID) [above right=2.5cm and 2.5cm of S] {ID/KW};
        \node[accept] (NUM) [above right=1.0cm and 2.5cm of S] {NUM};
        \node[state] (EQ) [right=1.1cm of S] {$S_{=}$};
        \node[accept] (EQEQ) [above right=0.5cm and 1.5cm of EQ] {==};
        \node[accept] (ASS) [below right=0.5cm and 1.5cm of EQ] {=};
        \node[state] (MN) [below right=2cm and 2.5cm of S] {$S_{-}$};
        \node[accept] (ARR) [above right=0.5cm and 1.5cm of MN] {->};
        \node[accept] (MINUS) [below right=0.5cm and 1.5cm of MN] {$-$};

        \path (S) edge node[above,sloped,font=\tiny] {字母/\_} (ID)
              (S) edge node[above,sloped,font=\tiny] {数字} (NUM)
              (S) edge node[above,sloped,font=\tiny] {=} (EQ)
              (EQ) edge node[below,sloped,font=\tiny] {other} (ASS)
              (EQ) edge node[above,sloped,font=\tiny] {=} (EQEQ)
              (S) edge node[below,sloped,font=\tiny] {$-$} (MN)
              (MN) edge node[below,sloped,font=\tiny] {other} (MINUS)
              (MN) edge node[above,sloped,font=\tiny] {>} (ARR);
    \end{tikzpicture}
    \end{center}
    \vfill
    \textbf{核心策略}：读取字符 $\to$ 根据首字符分派 $\to$ 必要时前瞻一个字符区分
\end{frame}

\begin{frame}{关键字 vs 标识符}
    \begin{block}{识别策略}
        先按标识符规则扫描完整单词，再查关键字表
    \end{block}
    \vfill
    \begin{columns}[T]
    \begin{column}{0.5\textwidth}
        \textbf{示例 1：\texttt{if123}}
        \begin{itemize}
            \item 首字符 \texttt{i}，进入标识符状态
            \item 继续读入 \texttt{f,1,2,3}
            \item 得到 \texttt{if123}
            \item 不在关键字表中 $\Rightarrow$ \textbf{标识符}
        \end{itemize}
    \end{column}
    \begin{column}{0.5\textwidth}
        \textbf{示例 2：\texttt{if=123}}
        \begin{itemize}
            \item 读入 \texttt{i,f}，下一个是 \texttt{=}（非字母数字）
            \item 得到 \texttt{if} $\Rightarrow$ \textbf{关键字}
            \item 然后 \texttt{=} $\Rightarrow$ \textbf{赋值号}
            \item 然后 \texttt{123} $\Rightarrow$ \textbf{整数}
        \end{itemize}
    \end{column}
    \end{columns}
\end{frame}

% ============================================================
\section{语法分析}

\begin{frame}{语法分析方法：递归下降}
    \begin{block}{方法简介}
        \textbf{递归下降分析}（Recursive Descent Parsing）是一种自顶向下的分析方法。\\
        为文法中每个非终结符编写一个解析函数，通过前瞻Token选择产生式。
    \end{block}
    \vfill
    \begin{columns}[T]
    \begin{column}{0.5\textwidth}
        \textbf{优点}
        \begin{itemize}
            \item 代码直观、易于调试
            \item 可利用上下文消歧
            \item 错误信息精确
        \end{itemize}
    \end{column}
    \begin{column}{0.5\textwidth}
        \textbf{处理步骤}
        \begin{enumerate}
            \item 消除左递归
            \item 确定各产生式的FIRST集
            \item 编写递归函数
        \end{enumerate}
    \end{column}
    \end{columns}
\end{frame}

\begin{frame}[fragile,shrink]{左递归消除}
    \textbf{原始文法}（以加减表达式为例）：
    \[
    \langle\text{加减表达式}\rangle \to \langle\text{加减表达式}\rangle\ \langle\text{加减运算符}\rangle\ \langle\text{项}\rangle \mid \langle\text{项}\rangle
    \]
    \vfill
    \textbf{消除后}（改写为迭代循环）：
\begin{lstlisting}[language=RustLike, basicstyle=\tiny\ttfamily]
fn parse_add_sub(&mut self) -> Result<Expr> {
    let mut left = self.parse_mul_div()?;
    loop {
        let op = match self.peek() {
            Plus => BinOp::Add,
            Minus => BinOp::Sub,
            _ => break,
        };
        self.advance();
        let right = self.parse_mul_div()?;
        left = Expr::BinaryOp(left, op, right);
    }
    Ok(left)
}
\end{lstlisting}
\end{frame}

\begin{frame}{运算符优先级}
    通过\textbf{函数调用层级}实现优先级：
    \begin{table}
    \centering
    \begin{tabular}{ccl}
    \toprule
    \textbf{层级} & \textbf{运算符} & \textbf{函数} \\
    \midrule
    1（低） & \texttt{..} & \texttt{parse\_range\_expr} \\
    2 & \texttt{== != < <= > >=} & \texttt{parse\_comparison} \\
    3 & \texttt{+ -} & \texttt{parse\_add\_sub} \\
    4 & \texttt{* /} & \texttt{parse\_mul\_div} \\
    5 & \texttt{* \& \&mut} (一元) & \texttt{parse\_unary} \\
    6 & \texttt{[] .N} (后缀) & \texttt{parse\_postfix} \\
    7（高） & 字面量/括号 & \texttt{parse\_primary} \\
    \bottomrule
    \end{tabular}
    \end{table}
    \vfill
    低优先级函数调用高优先级函数 $\Rightarrow$ 高优先级运算先绑定
\end{frame}

\begin{frame}{增强型语义分析与安全性}
    \begin{block}{核心目标}
        在语法分析过程中集成语义校验，提升工具的健壮性。
    \end{block}
    \vfill
    \begin{columns}[T]
    \begin{column}{0.5\textwidth}
        \textbf{可变性校验}
        \begin{itemize}
            \item \texttt{let mut} 严格声明
            \item 禁止对不可变变量赋值
            \item 支持解引用可变性校验
            \item 指向 \texttt{\&mut T} 的解引用可写
        \end{itemize}
    \end{column}
    \begin{column}{0.5\textwidth}
        \textbf{类型与迭代校验}
        \begin{itemize}
            \item \textbf{类型一致性}：赋值与声明时的类型必须完全匹配
            \item \textbf{迭代合法性}：\texttt{for} 循环仅支持 \texttt{Range} 或 \texttt{Array}
        \end{itemize}
    \end{column}
    \end{columns}
\end{frame}

\begin{frame}[shrink]{已实现的语法规则总览}
    \begin{columns}[T]
    \begin{column}{0.5\textwidth}
        \textbf{基础（必需）}
        \begin{itemize}
            \footnotesize
            \item 0.1 变量属性 (\texttt{mut}/空)
            \item 0.2 类型 (\texttt{i32})
            \item 0.3 左值 (标识符)
            \item 1.1 基础程序
            \item 1.2 语句/语句串
            \item 1.3 返回语句
            \item 1.4 函数输入（形参列表）
            \item 1.5 函数输出（返回类型）
            \item 2.0 变量声明
            \item 2.1 变量声明语句
            \item 2.2 赋值语句
            \item 3.1--3.5 表达式/运算/函数调用
            \item 4.1 if 选择
            \item 5.0--5.1 循环/while
        \end{itemize}
    \end{column}
    \begin{column}{0.5\textwidth}
        \textbf{扩展（额外实现）}
        \begin{itemize}
            \footnotesize
            \item 2.3 变量声明赋值
            \item 4.2--4.3 else / else if
            \item 5.2 for 循环
            \item 5.3 loop 循环
            \item 5.4 break / continue
            \item 6.1 不可变变量
            \item 6.2--6.3 不可变/可变引用
            \item 6.4 借用（解引用）
            \item 7.0--7.4 函数表达式块
            \item 8.1--8.3 数组
            \item 9.1--9.3 元组
        \end{itemize}
    \end{column}
    \end{columns}
\end{frame}

% ============================================================
\section{系统架构}

\begin{frame}{模块架构}
    \begin{center}
    \resizebox{\linewidth}{!}{
    \begin{tikzpicture}[
        box/.style={rectangle, draw, rounded corners, minimum width=2.5cm, minimum height=0.9cm, fill=blue!8, font=\small\bfseries},
        arrow/.style={->, >=Latex, thick}
    ]
        \node[box] (src) at (0,0) {源代码};
        \node[box] (lex) at (3.5,0) {Lexer};
        \node[box] (tok) at (7,0) {Token流};
        \node[box] (par) at (10.5,0) {Parser};
        \node[box] (ast) at (14,0) {AST};
        \node[box, fill=green!8] (tkm) at (3.5,-1.5) {token.rs};
        \node[box, fill=green!8] (asm) at (10.5,-1.5) {ast.rs};

        \draw[arrow] (src) -- (lex);
        \draw[arrow] (lex) -- (tok);
        \draw[arrow] (tok) -- (par);
        \draw[arrow] (par) -- (ast);
        \draw[arrow, dashed] (tkm) -- (lex);
        \draw[arrow, dashed] (asm) -- (par);
    \end{tikzpicture}
    }
    \end{center}
    \vfill
    \begin{columns}[T]
    \begin{column}{0.5\textwidth}
        \begin{itemize}
            \item \textbf{token.rs}：Token类型定义
            \item \textbf{lexer.rs}：词法扫描器
        \end{itemize}
    \end{column}
    \begin{column}{0.5\textwidth}
        \begin{itemize}
            \item \textbf{ast.rs}：AST节点 + 打印
            \item \textbf{parser.rs}：递归下降分析器
        \end{itemize}
    \end{column}
    \end{columns}
\end{frame}

% ============================================================
\section{测试与演示}

\begin{frame}[fragile,shrink]{测试用例：基础程序 (1.1--1.5)}
\begin{columns}[T]
\begin{column}{0.5\textwidth}
\begin{lstlisting}[language=RustLike, basicstyle=\tiny\ttfamily]
fn program_1_1() {
}
fn program_1_2() {
    ;;;;;;
}
fn program_1_3() {
    return ;
}
fn program_1_4(mut a:i32) {
}
fn program_1_5() -> i32 {
    return 1;
}
\end{lstlisting}
\end{column}
\begin{column}{0.5\textwidth}
\textbf{AST 输出：}
\begin{lstlisting}[basicstyle=\tiny\ttfamily,numbers=none]
FunctionDecl: program_1_1()
  Block

FunctionDecl: program_1_2()
  Block
    EmptyStmt (x6)

FunctionDecl: program_1_3()
  Block
    Return (void)

FunctionDecl: program_1_4(mut a: i32)
  Block

FunctionDecl: program_1_5() -> i32
  Block
    Return
      Int(1)
\end{lstlisting}
\end{column}
\end{columns}
\end{frame}

\begin{frame}[fragile]{测试用例：变量与赋值 (2.1--2.3)}
\begin{columns}[T]
\begin{column}{0.5\textwidth}
\begin{lstlisting}[language=RustLike, basicstyle=\tiny\ttfamily]
fn program_2_1() {
    let mut a;
    let mut b:i32;
}
fn program_2_2(mut a:i32) {
    a=32;
}
fn program_2_3() {
    let mut a=1;
    let mut b:i32=1;
}
\end{lstlisting}
\end{column}
\begin{column}{0.5\textwidth}
\textbf{AST 输出：}
\begin{lstlisting}[basicstyle=\tiny\ttfamily,numbers=none]
FunctionDecl: program_2_1()
  Block
    VarDecl: let mut a
    VarDecl: let mut b: i32

FunctionDecl: program_2_2(mut a: i32)
  Block
    Assign
      Ident(a)
      Int(32)

FunctionDecl: program_2_3()
  Block
    VarDeclAssign: let mut a = ...
      Int(1)
    VarDeclAssign: let mut b: i32 = ...
      Int(1)
\end{lstlisting}
\end{column}
\end{columns}
\end{frame}

\begin{frame}[fragile,shrink]{测试用例：控制流 (4.1--4.3, 5.1)}
\begin{columns}[T]
\begin{column}{0.5\textwidth}
\begin{lstlisting}[language=RustLike, basicstyle=\tiny\ttfamily]
fn program_4_3(a:i32) -> i32 {
    if a>0 {
        return a+1;
    } else if a<0 {
        return a-1;
    } else {
        return 0;
    }
}
fn program_5_1(mut n:i32) {
    while n>0 {
        n=n-1;
    }
}
\end{lstlisting}
\end{column}
\begin{column}{0.5\textwidth}
\textbf{AST 输出：}
\begin{lstlisting}[basicstyle=\tiny\ttfamily,numbers=none]
FunctionDecl: program_4_3(a: i32)
  Block
    If
      BinOp(>)
        Ident(a), Int(0)
      Block
        Return BinOp(+) ...
    ElseIf
      If
        BinOp(<)
          Ident(a), Int(0)
        ...
      Else
        Block
          Return Int(0)

FunctionDecl: program_5_1(mut n: i32)
  Block
    While
      BinOp(>): Ident(n), Int(0)
      Block
        Assign: n = n-1
\end{lstlisting}
\end{column}
\end{columns}
\end{frame}

\begin{frame}[fragile]{测试用例：数组与元组 (8.1--9.3)}
\begin{columns}[T]
\begin{column}{0.5\textwidth}
\begin{lstlisting}[language=RustLike, basicstyle=\tiny\ttfamily]
fn program_8_2(mut a:[i32;3]) {
    a=[1,2,3];
}
fn program_8_3(mut a:[i32;3]) {
    let mut b:i32=a[0];
    a[0]=1;
}
fn program_9_2(mut a:(i32,i32)) {
    a=(1,2);
}
fn program_9_3(mut a:(i32,i32)) {
    let mut b:i32=a.0;
    a.0=1;
}
\end{lstlisting}
\end{column}
\begin{column}{0.5\textwidth}
\textbf{AST 输出：}
\begin{lstlisting}[basicstyle=\tiny\ttfamily,numbers=none]
FunctionDecl: program_8_2(...)
  Block
    Assign
      Ident(a)
      ArrayLiteral(3 elems)
        Int(1), Int(2), Int(3)

FunctionDecl: program_8_3(...)
  Block
    VarDeclAssign: let mut b: i32
      Index: Ident(a), Int(0)
    Assign
      Index: Ident(a), Int(0)
      Int(1)

... (元组类似)
\end{lstlisting}
\end{column}
\end{columns}
\end{frame}

\begin{frame}{测试结果}
    \begin{alertblock}{全部测试通过}
        \textbf{28 个}测试函数覆盖规则 0.1 -- 9.3，全部成功通过词法分析和语法分析。
    \end{alertblock}
    \vfill
    \begin{block}{统计信息}
        \begin{itemize}
            \item 词法分析：正确识别所有关键字、标识符、数值、运算符、界符等
            \item 语法分析：成功构建完整AST，无报错
            \item 程序输出：``语法分析成功！''
        \end{itemize}
    \end{block}
\end{frame}

% ============================================================
\section{思考与总结}

\begin{frame}{思考：LR(1) 的局限}
    \begin{block}{问题}
        PPT提到``存在极个别的羁绊会使得LR(1)无法处理''
    \end{block}
    \vfill
    \textbf{典型案例：表达式块 vs 语句块}
    \begin{itemize}
        \item 遇到 \texttt{\{} 时，LR(1) 无法确定是语句块还是表达式块
        \item 递归下降可以通过\textbf{调用上下文}区分
        \item LR(1) 中可尝试改写产生式或使用 GLR 来处理
    \end{itemize}
    \vfill
    \textbf{元组与括号的歧义}
    \begin{itemize}
        \item \texttt{(expr)} 是括号表达式
        \item \texttt{(expr,)} 是单元素元组
        \item 需要看到逗号后才能区分 $\Rightarrow$ 需要更多前瞻
    \end{itemize}
\end{frame}

\section{环境与执行}

\begin{frame}{项目编译与执行}
    \begin{block}{编译环境}
        \begin{itemize}
            \item Rust 1.70+, Cargo
            \item LaTeX (XeLaTeX), latexmk (用于生成报告)
        \end{itemize}
    \end{block}
    \vfill
    \begin{block}{常用编译与执行命令}
        \begin{itemize}
            \item \texttt{cargo build --release} : 编译生成 release 二进制文件
            \item \texttt{cargo run --release} : 运行内置语法树解析测试
            \item \texttt{cargo run --release -- tests/test\_all.rs.txt} : 指定外部程序测试
        \end{itemize}
    \end{block}
\end{frame}

\begin{frame}{总结}
    \begin{block}{完成内容}
        \begin{itemize}
            \item 使用 \textbf{Rust} 实现了类Rust语言的词法分析器和语法分析器
            \item 词法分析采用手写C \textbf{ursor扫描器}，支持13种关键字和所有词法单元
            \item 语法分析采用\textbf{递归下降}方法，消除左递归，正确处理运算符优先级
            \item 实现了从基础规则（0.1--5.1）到扩展规则（5.2--9.3）的\textbf{全部34条}产生式
            \item 所有 28 个测试用例\textbf{全部通过}
        \end{itemize}
    \end{block}
    \vfill
    \begin{block}{个人收获}
        \begin{itemize}
            \item 深入理解了词法分析的状态机原理和语法分析的递归下降方法
            \item 掌握了左递归消除和运算符优先级的处理技巧
            \item 体会到 Rust 枚举 + 模式匹配在编译器前端实现中的优势
        \end{itemize}
    \end{block}
\end{frame}

\begin{frame}{参考资料}
    \begin{itemize}
        \item 陈火旺. 程序设计语言编译原理（第3版）. 国防工业出版社.
        \item Alfred V. Aho et al. Compilers: Principles, Techniques and Tools (Second Edition). 赵建华等译. 机械工业出版社.
        \item The Rust Reference \& Rust By Example.
    \end{itemize}
\end{frame}

\begin{frame}
    \begin{center}
        \Huge 谢谢！
    \end{center}
\end{frame}

\end{document}
